\section*{章末练习}

\begin{longtable}{L{0.08\linewidth}L{0.32\linewidth}L{0.5\linewidth}}
\toprule
题号 & 类型 & 题目 \\
\midrule
1 & 时延组成 & 将机械硬盘随机读时延拆成寻道、旋转等待、传输和队列四部分；顺序读主要减少哪几项？ \\
2 & 调度 & 比较 FCFS 与电梯式调度对吞吐、公平性和尾延迟的影响。 \\
3 & 扇区格式 & 用户数据之外为何还需要同步、地址标记与 ECC？这些开销为何不会出现在标称逻辑容量中？ \\
4 & SATA 路径 & 从主机命令到盘片读取再到 DMA 完成，列出 SATA 控制器与磁盘内部的主要阶段。 \\
5 & 写缓存 & 开启易失写缓存后，命令完成能说明什么？文件系统需要何种接口建立持久化顺序？ \\
6 & 故障征兆 & 重映射扇区数、待定扇区、CRC 错误和命令超时分别更可能指向哪些故障域？ \\
\bottomrule
\end{longtable}

\section*{参考解答与要点}

\begin{longtable}{L{0.08\linewidth}L{0.32\linewidth}L{0.5\linewidth}}
\toprule
题号 & 类型 & 解答要点 \\
\midrule
1 & 时延组成 & 随机读包含队列等待、磁头寻道、等目标扇区转到磁头下和数据传输；顺序读主要摊薄寻道与旋转等待。 \\
2 & 调度 & FCFS 顺序公平且简单，但磁头移动大；电梯式重排可提高局部性和吞吐，却需等待上界或优先级机制防止远端请求长期延迟。 \\
3 & 扇区格式 & 同步用于找边界，地址标记确认位置，ECC 检测纠正介质位错；它们属于物理格式保护开销，逻辑块接口只暴露用户数据容量。 \\
4 & SATA 路径 & 主机准备命令/缓冲，HBA 发 FIS；磁盘排队并寻道、等待旋转、读扇区与 ECC；数据经链路到 HBA，DMA 入内存，状态完成并通知 CPU。 \\
5 & 写缓存 & 可能只表示数据进入盘内易失缓存；上层需用 FLUSH/FUA 等已协商语义，并正确安排数据与元数据命令顺序。 \\
6 & 故障征兆 & 重映射/待定扇区偏向介质表面或磁头路径，CRC 错误偏向线缆/连接/信号，超时还可能来自供电、控制器、固件或设备内部停滞，需结合多项证据判断。 \\
\bottomrule
\end{longtable}
